\chapter{Metodología del Proyecto} \label{cap:metodologia}

El proyecto se desarrolló bajo el marco de trabajo Scrum, estructurado en ciclos iterativos (sprints) de dos semanas que permitieron el diseño y calibración adaptativa del banco de pruebas de red. Las actividades se organizaron a partir de un Product Backlog trazado con los objetivos del trabajo, validando de forma incremental tanto la infraestructura automatizada del testbed como el algoritmo del recomendador interactivo.

\section{Product Backlog y trazabilidad con los objetivos}

El \emph{Product Backlog} se derivó directamente de los cuatro objetivos específicos del trabajo, de modo que cada conjunto de tareas contribuye de forma trazable a un objetivo. La \cref{tab:backlog-oe} resume esta correspondencia.

\begin{table}[htbp]
\centering
\scriptsize
\renewcommand{\arraystretch}{1.25}
\caption{Mapeo del Product Backlog con los objetivos específicos (OE).}
\label{tab:backlog-oe}
\begin{tabularx}{\textwidth}{|c|Y|}
\hline
\textbf{Objetivo} & \textbf{Líneas de trabajo del backlog} \\
\hline
OE1 & Testbed automatizado (IaC), \emph{baselines} de rendimiento de los cuatro CNI, rightsizing de recursos. \\
\hline
OE2 & Biblioteca de Network Policies, validación automatizada, observabilidad y alertas de red. \\
\hline
OE3 & Matriz de criterios y métricas, normalización, ponderaciones y modelo MCDA. \\
\hline
OE4 & Guía metodológica, prototipo recomendador (SPA) y empaquetado de entregables. \\
\hline
\end{tabularx}
\end{table}

\section{Planificación de Sprints (Sprint Planning)}

El proyecto se planificó en \textbf{17 sprints} (Sprint 0 a Sprint 16) de dos semanas cada uno, entre el 1.\textsuperscript{o} de julio de 2025 y el 19 de febrero de 2026, según el cronograma del anteproyecto~\cite{ref29}. La \cref{tab:sprints} presenta el \emph{Sprint Planning} maestro, donde cada sprint se asocia a su entregable principal y al objetivo que atiende.

\begin{table}[htbp]
\centering
\scriptsize
\renewcommand{\arraystretch}{1.2}
\caption{Planificación maestra de sprints y su relación con los objetivos específicos.}
\label{tab:sprints}
\begin{tabularx}{\textwidth}{|c|Y|c|c|}
\hline
\textbf{Sprint} & \textbf{Foco / Entregable} & \textbf{Fechas} & \textbf{OE} \\
\hline
0  & Kickoff, entorno y tablero & 01--10/07/2025 & --- \\
\hline
1  & Kit de laboratorio v0.1 & 11--24/07/2025 & OE1 \\
\hline
2  & CNI \#1 \emph{baseline} & 25/07--07/08/2025 & OE1 \\
\hline
3  & Biblioteca de Network Policies v0.1 & 08--21/08/2025 & OE2 \\
\hline
4  & CNI \#2 \emph{baseline} & 22/08--04/09/2025 & OE1 \\
\hline
5  & Matriz de criterios v0.5 + checklist v0.1 & 05--18/09/2025 & OE3 \\
\hline
6  & Validación automatizada de Network Policies & 19/09--02/10/2025 & OE2 \\
\hline
7  & CNI \#3 \emph{baseline} & 03--16/10/2025 & OE1 \\
\hline
8  & Guía metodológica v0.6 & 17--30/10/2025 & OE3/OE4 \\
\hline
9  & Observabilidad y alertas de red & 31/10--13/11/2025 & OE1/OE2 \\
\hline
10 & CNI \#4 \emph{baseline} o variante & 14--27/11/2025 & OE1 \\
\hline
11 & \emph{Rightsizing} v0.1 & 28/11--11/12/2025 & OE1 \\
\hline
12 & Biblioteca de NP v1.0 + \emph{snippets} & 12--25/12/2025 & OE2 \\
\hline
13 & Matriz v1.0 + checklist v1.0 & 26/12/2025--08/01/2026 & OE3 \\
\hline
14 & Guía metodológica v1.0 + empaquetado & 09--22/01/2026 & OE4 \\
\hline
15 & Integración final y reporte & 23/01--05/02/2026 & --- \\
\hline
16 & Pulido final y entrega v1.0 & 06--19/02/2026 & --- \\
\hline
\end{tabularx}
\end{table}

El diseño técnico que materializa estos sprints se presenta en los \cref{sec:Metodos,cap:herramienta}. La \cref{sec:desarrollo} detalla la ejecución iterativa, agrupada por fases temáticas y con la retrospectiva de cada agrupación.

\section{Desarrollo iterativo por fases} \label{sec:desarrollo}

El desarrollo se ejecutó a lo largo de los diecisiete sprints planificados (\cref{tab:sprints}). Para facilitar la lectura, esta sección agrupa los sprints en fases temáticas y cierra cada fase con una retrospectiva que sintetiza los aprendizajes ---el artefacto de mejora continua propio de Scrum---, en lugar de narrar sprint por sprint. El detalle técnico de cada componente se encuentra en los \cref{sec:Metodos,cap:herramienta}.

\subsection{Fase de preparación del laboratorio (Sprints 0--1)}

Los primeros sprints establecieron las bases del proyecto: repositorios, integración continua, tablero de tareas y el \emph{kit de laboratorio} v0.1. Se desplegó la pila de observabilidad (Prometheus y Grafana) y se construyeron los primeros \emph{scripts} de inyección de tráfico (iperf3) con su validación intra e inter-nodo.

\textbf{Retrospectiva.} La automatización temprana del entorno como infraestructura como código resultó decisiva: garantizó que cada CNI se evaluara en condiciones idénticas. Se ajustó la frecuencia de \emph{scraping} de métricas para capturar los picos de las ráfagas de tráfico sin introducir carga adicional.

\subsection{Fase de evaluación de rendimiento de los CNI (Sprints 2, 4, 7 y 10)}

Cada CNI ---Flannel, Calico, Cilium y Antrea--- se evaluó en su propio sprint de \emph{baseline}: instalación, pruebas de throughput y latencia (percentiles p95/p99), Network Policies mínimas y un análisis corto. El protocolo de rotación reemplazaba por completo la infraestructura entre ciclos, partiendo siempre de nodos limpios.

\textbf{Retrospectiva.} Separar cada CNI en un sprint independiente evitó la contaminación cruzada de configuraciones. El principal impedimento ---el desajuste de MTU al instalar Cilium--- se resolvió calculando la MTU efectiva y se documentó como aprendizaje reutilizable para los CNI restantes.

\subsection{Fase de seguridad y micro-segmentación (Sprints 3, 6 y 12)}

Esta fase construyó la biblioteca de Network Policies bajo el enfoque Zero Trust: casos de \emph{Default Deny}, segmentación por capas \texttt{frontend $\rightarrow$ backend $\rightarrow$ database} y reglas avanzadas (L7 en Cilium, \emph{Tiers} en Antrea). Se implementó además el pipeline de validación automatizada que verifica, para cada par de Pods, si la conexión debe permitirse o bloquearse.

\textbf{Retrospectiva.} La validación automatizada confirmó que la sola declaración de una política no garantiza su \emph{enforcement}: Flannel aceptó los objetos pero no los aplicó. Este hallazgo telemático ---la diferencia entre declarar y aplicar una política--- se convirtió en una de las variables centrales del modelo de decisión.

\subsection{Fase del modelo de decisión MCDA (Sprints 5 y 13)}

Se definieron los criterios y métricas de evaluación, los factores de ponderación por perfil y la restricción no compensatoria de seguridad, consolidados en la matriz de criterios v1.0 y el \emph{checklist} de configuración. El modelo se ensayó primero con datos de prueba y luego con los datos reales del \emph{testbed}.

\textbf{Retrospectiva.} Parametrizar los pesos por perfil ---en lugar de fijar una única recomendación--- hizo el modelo adaptable a distintas organizaciones. El ensayo con datos sintéticos antes de los reales permitió detectar y corregir el sesgo de magnitud entre métricas con unidades distintas.

\subsection{Fase de la herramienta y la guía (Sprints 8 y 14)}

Estos sprints produjeron la guía metodológica (de la v0.6 a la v1.0) y el empaquetado de los entregables, incluyendo el prototipo recomendador. Se redactaron los procedimientos paso a paso para instalar y probar cada CNI, base de los manuales de usuario e instalación incluidos como anexos.

\textbf{Retrospectiva.} Documentar la guía en paralelo al desarrollo ---y no al final--- evitó la pérdida de conocimiento por la complejidad de las configuraciones. La estructura ``cómo se hace'' demostró ser más útil para equipos sin especialización que una descripción puramente teórica.

\subsection{Fase de observabilidad y \emph{rightsizing} (Sprints 9 y 11)}

Se ampliaron las métricas de red (L3--L7), los \emph{dashboards} y las reglas de alerta, y se desarrolló el método de \emph{rightsizing}: relacionar el consumo de CPU y memoria del agente CNI con el throughput para dimensionar la infraestructura sin márgenes precautorios excesivos.

\textbf{Retrospectiva.} La correlación entre consumo de recursos y rendimiento cuantificó el costo telemático real de cada CNI, base del aporte sobre reducción de sobredimensionamiento y su impacto económico en las organizaciones.

\subsection{Fase de integración y cierre (Sprints 15--16)}

Los últimos sprints integraron todos los componentes, ejecutaron las pruebas finales y consolidaron resultados, cerrando con el control de calidad documental y de laboratorio y la entrega v1.0.

\textbf{Retrospectiva.} La integración incremental a lo largo del proyecto ---y no en una única fase final--- redujo el riesgo de incompatibilidades. El control de calidad final se concentró en la consistencia entre los datos del \emph{testbed}, el modelo y el prototipo.
